\documentclass[10.5pt]{article}
\usepackage[margin=0.74in]{geometry}
\usepackage{booktabs,graphicx,hyperref,xcolor,enumitem,tabularx,microtype,fontspec}
\setmainfont{TeX Gyre Heros}
\hypersetup{colorlinks=true,urlcolor=blue!55!black,linkcolor=blue!55!black}
\setlist{nosep,leftmargin=1.35em}
\setlength{\parskip}{0.42em}
\setlength{\parindent}{0pt}
\newcommand{\code}[1]{\texttt{#1}}
\newcommand{\figdir}{artifacts/reports/qwen-code-replacement-study-20260915/figures}
\title{Native Qwen Code SFT Replacement Study\\\large An end-to-end pilot with an independent TRExFitter scorer}
\author{ATLAS H$\rightarrow\gamma\gamma$ coding-agent project}
\date{15 September 2026}

\begin{document}
\maketitle

\begin{abstract}
This pilot replaces the earlier repository-specific Qwen tool loop with Qwen
Code 0.23.4 itself.  Canonical training rows remain structured conversations;
the Qwen tokenizer renders its own function-call tokens.  A Qwen3.5-4B model
was trained for one epoch with LoRA on four A100 GPUs and evaluated beside the
unchanged base checkpoint on the same six held-out tasks.  Base scored 3/6 and
SFT scored 2/6 under the strict external gate.  Thus this run validates the
replacement plumbing but does not establish an SFT accuracy improvement.
\end{abstract}

\section{Question and workflow}
The practical question is whether a general coding model can translate a
physicist's request into a small, correct change to \code{analysis.config},
preserve unrelated analysis choices, run the specified Python verifier, and
truthfully report the observed result.  The model is a reasoning-capable Qwen
checkpoint, but hidden chain-of-thought is neither stored nor supervised:
training and data rendering consistently set \code{enable\_thinking=False}.

The replacement protocol is:
\begin{enumerate}
  \item Qwen Code exposes exactly \code{read\_file}, \code{edit}, and
        \code{run\_shell\_command}; the JSON schemas are captured from the
        official \code{@qwen-code/qwen-code} 0.23.4 request and pinned.
  \item Qwen Code parses calls, applies edits in a disposable task workspace,
        runs shell commands, and returns native tool results to the model.
  \item An independent repository scorer inspects the final file and event
        stream.  It is not available to the model as a special-purpose tool.
\end{enumerate}
The short study system prompt restricts work to the task workspace and asks for
minimal edits plus the exact verifier command.  The same prompt, decoding
endpoint, six tasks, 4,096-token output limit, and scorer are used for base and
SFT.  Qwen Code runs in \code{yolo} approval mode only because every task uses
an isolated disposable copy.

\section{Dataset and rendering}
The refreshed public release has 130 train and 39 validation conversations:
156 replay-approved deterministic demonstrations, 11 real Codex/OpenCode
teacher trajectories, and two externally accepted Qwen Code trajectories.
The Qwen additions are one clean validation success and one training recovery
(failed no-side-effect calls followed by a correct edit and verifier run).
They were collected after this checkpoint was trained and therefore are not
part of its supervision.  The evaluated checkpoint used the immediately prior
129-train/38-validation Qwen Code-native snapshot; the two new native rows are
for the next training round.
Records store structured \code{messages}, optional \code{tools}, provenance,
and verification evidence---not rendered markup, hidden reasoning, token IDs,
or labels.  The target checkpoint's chat template creates model-specific
tokens at training time.  Preflight proved identical full and turn-wise
rendering, assistant-only loss masks, no-thinking behavior, and no truncation
at the 16,384-token context length.

\section{Training}
VERL's SFT trainer ran Qwen3.5-4B for one epoch: LoRA rank 16 and alpha 16,
learning rate $10^{-4}$, BF16, global batch 8, seed 1, four A100 GPUs, and 16
optimizer steps.  LoRA trains small low-rank updates while leaving base weights
frozen.  Peak allocated GPU memory was 24.65~GB per rank (33.74~GB reserved).
The adapter was reconstructed from VERL's sharded checkpoint and merged only
for native serving.

\begin{figure}[ht]
\centering\includegraphics[width=.82\textwidth]{\figdir/training-loss.pdf}
\caption{Prediction loss over the single training epoch.  The endpoint is low,
but loss measures imitation of held-out conversations, not successful execution.}
\end{figure}

Training loss fell from 0.375 to 0.056; held-out assistant-token loss was
0.079.  These values show that the model learned the demonstration distribution.
They do not prove that it learned robust tool selection, editing, and recovery.

\section{Matched result}
\begin{figure}[ht]
\centering\includegraphics[width=.72\textwidth]{\figdir/matched-strict-success.pdf}
\caption{Strict end-to-end success on the same six held-out tasks.  Six tasks
are a smoke test, not a statistically precise benchmark.}
\end{figure}

Base Qwen3.5-4B passed 3/6 (50\%); the one-epoch LoRA checkpoint passed 2/6
(33\%).  Both solved the held-out MINOS repair and ZH filename repair.  Base
also correctly recognized the ggH-directory no-op.  This pilot therefore gives
no evidence that the current SFT improves strict accuracy.  The difference is
one task and is far too small for a stable effect estimate.  The most useful
result is diagnostic: low validation loss coexists with end-to-end regressions.
A larger held-out evaluation and training directly on more Qwen Code-native
successful and recovery traces are required before tuning scale or epochs.

\section{What ``strict success'' means}
A task passes only when all gates pass:
\begin{itemize}
  \item \textbf{requested state}: the named setting has the requested semantic value;
  \item \textbf{preservation}: unrelated settings were not changed;
  \item \textbf{static validity}: the Python parser and requested evidence checks pass;
  \item \textbf{observed verification}: after its final edit, the agent successfully
        ran the exact requested verifier command;
  \item \textbf{native scope}: only allowed native tools and workspace-local paths were used.
\end{itemize}
A plausible final message or a parseable file alone is insufficient.  Static
validity is the SFT-stage target.  Running TRExFitter and comparing output
histograms is a stronger semantic-equivalence reward intended for later RL.

\section{Diagnostic gates and examples}
\begin{tabularx}{\textwidth}{@{}p{.18\textwidth}p{.26\textwidth}X@{}}
\toprule
Failure & Meaning & Concrete matched-run example \\
\midrule
Editing / contract & The tools ran, but the requested semantic value was wrong,
or unrelated analysis choices changed. & Base: scan-steps task left
\code{doLHscan} wrong. SFT: MINOS task wrote the wrong \code{UseMinos};
scan-steps also changed \code{FitType}. \\
Protocol & The agent did not complete the required edit--verify workflow. &
Base attempted unavailable \code{grep\_search} on the MINOS task and never
completed verification. SFT made edits but did not obtain a successful exact
verifier call on the CPU-repair task. \\
Verifier / evidence & The exact verifier did not succeed after the last edit,
or its requested evidence level failed. & Base MINOS ended after an API error
before verification. SFT ggH no-op damaged available-directory settings, so
static validation and evidence failed. \\
Scope & A successful command or path escaped the disposable task workspace, or
an unapproved native tool was requested. & Base CPU repair attempted
\code{grep\_search} and an outside path. SFT CPU repair referenced a path outside
the task workspace. \\
Harness / API & The native CLI or model endpoint ended abnormally; this is kept
separate from edit correctness. & Base MINOS returned an API 400 after editing.
SFT ggH no-op ended with a CLI error after its incorrect edits. \\
\bottomrule
\end{tabularx}

Failures can overlap: for example, a wrong edit can make validation fail, and a
subsequent out-of-scope command can also fail the protocol gate.  Keeping the
gates separate prevents an editing regression from being mislabeled as mere
tool syntax trouble.

\section{Successful and recovery examples}
The clean Qwen validation example reads the Fit block, uncomments
\code{UseMinos: mu\_H}, runs the exact verifier, observes success, and reports
only that evidence.  The admitted training recovery first makes two failed
path reads (no side effects), then reads the workspace file, repairs
\code{FitRegion}, and observes a successful verifier result.  Recovery traces
are valuable because a coding agent must respond to real tool errors; they are
admitted only when the final state passes the independent scorer.

\section{Conclusions and next experiment}
The replacement architecture is now credible: exact native Qwen Code schemas,
real native execution, model-neutral storage, Qwen-owned rendering, and an
independent domain scorer.  The first 4B LoRA result is negative on strict
accuracy despite a good loss curve.  The next defensible experiment is to add
more independently passing native Qwen successes and recoveries, reserve a
larger untouched test set, and compare base/SFT with multiple decoding seeds.
Only then should we attribute changes to LoRA rank, learning rate, model size,
or additional epochs.

\section*{Reproducibility pointers}
Dataset: \href{https://huggingface.co/datasets/cxyang-ucb/hyy-sft}{cxyang-ucb/hyy-sft}.
The repository records the Qwen Code 0.23.4 tool snapshot, canonical split
hashes, VERL checkpoint metrics, native event streams, exact task workspaces,
and independent per-task scores.  Checkpoints and raw event streams are kept in
local artifacts and are not part of the public dataset.

\end{document}
